\begin{table}[t]
\centering
\caption{Standardized Effect Sizes}
\label{tab:sde}
\begin{tabular}{lcccccc}
\hline\hline
Outcome & $\hat{\beta}$ & SE & SD($Y$) & SDE & SE(SDE) & Classification \\
\hline
\multicolumn{7}{l}{\textit{Panel A: Pooled}} \\
Establishments & 0.0020 & 0.0048 & 1.0959 & 0.0255 & 0.0623 & Small positive \\
Population & -0.0012 & 0.0039 & 1.1342 & -0.0145 & 0.0488 & Small negative \\
Steuerkraft & -0.0305 & 0.0090 & 1.3575 & -0.3191 & 0.0946 & Large negative \\
\hline
\multicolumn{7}{l}{\textit{Panel B: Steuerkraft by Municipality Size}} \\
Small municipalities & -0.0291 & 0.0180 & 0.8084 & -0.5118 & 0.3168 & Large negative \\
Large municipalities & -0.0299 & 0.0099 & 0.9124 & -0.4661 & 0.1545 & Large negative \\
\hline\hline
\end{tabular}
\begin{tablenotes}
\item \textit{Notes:} \textbf{Country:} Switzerland. \textbf{Research question:} Does the corporate--personal tax multiplier wedge in Swiss municipalities cause factor-specific sorting of firms and residents, or does it shift the tax base without physical relocation? \textbf{Policy mechanism:} Swiss municipalities in the Canton of Zurich set separate Steuerfuss (tax multiplier) rates for natural persons and legal persons (corporations), creating within-municipality variation in relative tax prices faced by different economic agents. \textbf{Outcome definition:} Panel A outcomes are log establishments (STATENT count of business locations), log population (permanent residents), and log Steuerkraft (3-year rolling average municipal tax capacity in CHF millions). Panel B splits Steuerkraft by municipality size. \textbf{Treatment:} Continuous; corporate Steuerfuss rate in percentage points (mean 119, SD 14.2). \textbf{Data:} Canton of Zurich Statistical Office (Steuerfuss, Steuerkraft), Federal Statistical Office STATENT (establishments), BFS population statistics, 2012--2023, municipality-year panel. \textbf{Method:} OLS with municipality and year fixed effects, standard errors clustered at municipality level. \textbf{Sample:} 172 municipalities in the Canton of Zurich with complete data on separate Steuerfuss rates for natural and legal persons. SDE $= \hat{\beta} \times \text{SD}(X) / \text{SD}(Y)$ where SD($X$) is the standard deviation of the corporate Steuerfuss and SD($Y$) is the pre-treatment (2012--2014) standard deviation of the outcome. Classification refers to magnitude, not statistical significance: Large ($|$SDE$| > 0.15$), Moderate ($0.05$--$0.15$), Small ($0.005$--$0.05$), Null ($< 0.005$).
\end{tablenotes}
\end{table}
